\documentclass[12pt,a4paper]{article}
% !TeX program = xelatex
\usepackage{fontspec}
\setmainfont{DejaVu Serif}
\usepackage[russian]{babel}
\usepackage{geometry}
\geometry{left=20mm,right=20mm,top=20mm,bottom=22mm}
\usepackage{amsmath,amssymb,mathtools}
\usepackage{array,booktabs,longtable}
\usepackage{xcolor}
\usepackage{hyperref}
\usepackage{tikz}
\hypersetup{colorlinks=true,linkcolor=blue!50!black,urlcolor=blue!50!black}
\setlength{\parindent}{0pt}
\setlength{\parskip}{5pt}
\setlength{\emergencystretch}{2em}
\renewcommand{\arraystretch}{1.25}
\newcommand{\err}{\textcolor{red}{Ошибка:}}
\newcommand{\idea}{\textcolor{blue!60!black}{Ключевая идея:}}
\newcommand{\result}{\textcolor{teal!60!black}{Итог:}}
\begin{document}
\thispagestyle{empty}
\vspace*{45mm}
\begin{center}
{\LARGE\bfseries Ревью домашней работы\par}
\vspace{22mm}
{\large Автор: Лёвин Артём Александрович\par}
\vspace{9mm}
{\large 21 октября 2026 года\par}
\end{center}
\vfill
\begin{center}
\begin{tikzpicture}[scale=1]
\draw[blue!45!black,line width=.7pt] (-5,0) .. controls (-2,0.65) and (2,-0.65) .. (5,0);
\end{tikzpicture}
\end{center}
\newpage
\section*{Сводная таблица}
\small
\begin{longtable}{@{}p{1.25cm}p{2.55cm}p{6.0cm}p{2.7cm}p{2.65cm}@{}}
\toprule
\textbf{№} & \textbf{Тема} & \textbf{Суть ошибки} & \textbf{Ваш ответ} & \textbf{Правильный ответ}\\
\midrule
\endhead
\hyperref[task:2]{2} & Сложные проценты & При переносе остатка на начало второго года записано 1626 вместо 1620 рублей; дальнейшие вычисления и окончательный ответ верны. & 1749,6 руб. & 1749,6 руб.\\
\midrule
\hyperref[task:8]{8} & Проценты и снятия & Во втором году после начисления процентов ошибочно получено 77,4 вместо 77,7 тысячи рублей; ошибка переносится на третий год. & 70\,770 руб. & 71\,085 руб.\\
\midrule
\hyperref[task:9]{9} & Обратная задача на вклад & Вместо множителя 1,2 для ставки 20\% использован множитель 1,1; затем коэффициент при неизвестной сумме записан непоследовательно. & $\approx114\,550$ руб. & 85\,000 руб.\\
\bottomrule
\end{longtable}
\normalsize
\newpage
\thispagestyle{empty}
\vspace*{\fill}
\begin{center}{\Huge\bfseries Разбор ошибок}\end{center}
\vspace*{\fill}
\newpage
\phantomsection\label{task:2}
\section*{Задание № 2}
\textbf{Текст задания.} На вклад поместили 1500 руб. под 8\% годовых на два года. Какова сумма в конце второго года, если пополнений и снятий нет?

\textbf{Ваше решение.} В таблице для первого года указаны начальная сумма 1500 руб., проценты 120 руб. и конечная сумма 1620 руб. Во втором году в столбце начальной суммы записано 1626 руб.; далее указаны проценты 129,6 руб. и конечная сумма 1749,6 руб.

\textbf{Ваш ответ.} 1749,6 руб.

\err\ Неверно перенесён остаток первого года в начало второго года. Конечный ответ при этом совпал с правильным.

\textbf{Где именно возникает ошибка.} Последняя правильная строка первого года:
\[
1500+1500\cdot0{,}08=1500+120=1620.
\]
Следовательно, при отсутствии операций со вкладом начало второго года должно равняться 1620 руб., а не записанным 1626 руб.

\textbf{Почему этот шаг неверен.} При сложных процентах накопленный за предыдущий год остаток без изменения становится суммой, на которую начисляют проценты в следующем году. Между первым и вторым годами по условию нет ни пополнения, ни снятия. Поэтому число 1626 руб. не может появиться: оно на 6 руб. больше полученных 1620 руб. Более того, указанные Вами проценты 129,6 руб. составляют 8\% именно от 1620 руб. Если бы начальная сумма второго года действительно была 1626 руб., проценты были бы равны $1626\cdot0{,}08=130{,}08$ руб., а не 129,6 руб. Значит, ошибка локальная: в одной ячейке неверно записано число, но следующий расчёт фактически выполнен от правильной суммы.

\idea\ Конечная сумма года без пополнений и снятий переходит в следующий год как начальная; проценты за каждый год начисляют на эту, уже увеличенную, сумму.

\textbf{Исправление от последнего правильного шага.} Заменим во втором году 1626 на 1620. Тогда проценты составят $1620\cdot0{,}08=129{,}6$ руб., а остаток --- $1620+129{,}6=1749{,}6$ руб.

\textbf{Полное правильное решение.} В первый год начислим восемь процентов на первоначальные 1500 руб.:
\[
P_1=1500\cdot\frac{8}{100}=120\ \text{руб.},\qquad
S_1=1500+120=1620\ \text{руб.}
\]
Так как в начале второго года операций со вкладом нет, начальная сумма равна 1620 руб. Начислим проценты повторно:
\[
P_2=1620\cdot\frac{8}{100}=129{,}6\ \text{руб.},\qquad
S_2=1620+129{,}6=1749{,}6\ \text{руб.}
\]
\textbf{Проверка.} По формуле сложных процентов:
\[
1500\cdot(1+0{,}08)^2=1500\cdot1{,}1664=1749{,}6.
\]
Оба способа дают одну сумму.

\result\ Ваш окончательный ответ верен: 1749,6 руб. Требуется исправить только запись начальной суммы второго года.

\textbf{Задача на закрепление.} В банк внесли 1200 руб. под 5\% годовых на два года без пополнений и снятий. Составьте таблицу с начальной суммой, начисленными процентами и конечной суммой за каждый год. Найдите остаток в конце второго года.
\newpage
\phantomsection\label{task:8}
\section*{Задание № 8}
\textbf{Текст задания.} На вкладе первоначально было 80 тыс. руб. Ставка --- 5\% годовых, срок три года. В начале второго и третьего годов снимают по 10 тыс. руб. Сколько рублей останется в конце срока?

\textbf{Ваше решение.} В первом году получены 4 тыс. руб. процентов и остаток 84 тыс. руб. Во втором году после снятия 10 тыс. руб. записана сумма 74 тыс. руб.; проценты 3,7 тыс. руб. вычислены верно, но остаток указан как 77,4 тыс. руб. В третьем году использована сумма 67,4 тыс. руб., начислены 3,37 тыс. руб. и получены 70,77 тыс. руб.

\textbf{Ваш ответ.} 70\,770 руб.

\err\ При сложении суммы вклада и процентов во втором году получено 77,4 вместо 77,7 тыс. руб.; поэтому остаток третьего года рассчитан от заниженной суммы.

\textbf{Где именно возникает ошибка.} До начисления процентов второго года всё верно:
\[
84-10=74,\qquad 74\cdot0{,}05=3{,}7
\quad\text{(тыс. руб.).}
\]
Первый неверный переход --- запись конечной суммы второго года:
\[
74+3{,}7=\boxed{77{,}7},\qquad\text{а не }77{,}4.
\]
\textbf{Почему этот шаг неверен.} Процентный доход не заменяет остаток на счёте: он \emph{прибавляется} к сумме, находившейся на вкладе после снятия. Здесь семь десятых прибавляются к семидесяти четырём целым, поэтому получается 77,7. Разница в 0,3 тыс. руб. появилась именно в сложении, а не в формуле процентов. В начале третьего года нужно вычесть 10 тыс. руб. из правильного остатка 77,7 тыс. руб. Если перенести неверные 77,4 тыс. руб., ошибка сохранится и дополнительно повлияет на проценты третьего года.

\idea\ Сначала выполняют снятие в начале года, затем на оставшуюся сумму начисляют годовые проценты и прибавляют их в конце года. В каждом следующем году используют точный остаток предыдущего.

\textbf{Исправление от последнего правильного шага.} После начисления 3,7 тыс. руб. во втором году:
\[
S_2=74+3{,}7=77{,}7\ \text{тыс. руб.}
\]
В начале третьего года:
\[
77{,}7-10=67{,}7\ \text{тыс. руб.}
\]
Проценты за третий год:
\[
67{,}7\cdot0{,}05=3{,}385\ \text{тыс. руб.}
\]
Итоговый остаток:
\[
67{,}7+3{,}385=71{,}085\ \text{тыс. руб.}
\]
\textbf{Полное правильное решение.} Будем считать в тысячах рублей, а в конце переведём сумму в рубли.

\textit{Первый год.} Доход равен $80\cdot0{,}05=4$ тыс. руб. Сумма на конец года $80+4=84$ тыс. руб.

\textit{Второй год.} В начале года снимают 10 тыс. руб., остаётся $84-10=74$ тыс. руб. Доход равен $74\cdot0{,}05=3{,}7$ тыс. руб. На конец года имеем $74+3{,}7=77{,}7$ тыс. руб.

\textit{Третий год.} Снова снимают 10 тыс. руб., остаётся $77{,}7-10=67{,}7$ тыс. руб. Доход равен $67{,}7\cdot0{,}05=3{,}385$ тыс. руб. На конец года остаётся $67{,}7+3{,}385=71{,}085$ тыс. руб.

Перевод единиц:
\[
71{,}085\cdot1000=71\,085\ \text{руб.}
\]
\textbf{Проверка.} Независимо пересчитаем три года с множителем 1,05:
\[
\bigl((80\cdot1{,}05-10)\cdot1{,}05-10\bigr)\cdot1{,}05
=71{,}085\ \text{тыс. руб.}
\]
Разница с Вашим ответом: $71\,085-70\,770=315$ руб.

\result\ Правильный остаток --- 71\,085 руб. Первая ошибка возникла при сложении в конце второго года.

\textbf{Задача на закрепление.} На вкладе было 60 тыс. руб. под 10\% годовых на три года. В начале второго и третьего годов снимают по 5 тыс. руб. Найдите остаток в конце третьего года в рублях.
\newpage
\phantomsection\label{task:9}
\section*{Задание № 9}
\textbf{Текст задания.} Вклад открыли на три года под 20\% годовых. В начале второго и третьего годов добавляли по 10 тыс. руб. В конце третьего года на счёте оказалось 173,28 тыс. руб. Найдите первоначальную сумму в тысячах рублей.

\textbf{Ваше решение.} Записана правильная по структуре модель $173{,}28=Sk^3+10k^2+10k$, где $S$ --- первоначальный вклад в тысячах рублей, $k$ --- годовой множитель. Далее вместо ставки 20\% подставлен множитель 1,1, соответствующий ставке 10\%. В последующих преобразованиях коэффициент при $S$ записан как 1,311, а ответ получен приблизительно 114,55 тыс. руб.

\textbf{Ваш ответ.} Приблизительно 114\,550 руб. (114,55 тыс. руб.).

\err\ Ставке 20\% поставлен в соответствие неверный множитель 1,1; кроме того, при переписывании коэффициента при неизвестной сумме допущена числовая несогласованность.

\textbf{Где именно возникает ошибка.} Последний содержательно верный шаг:
\[
173{,}28=Sk^3+10k^2+10k.
\]
Первый ошибочный шаг --- переход к $k=1{,}1$. По условию ставка равна 20\%, поэтому:
\[
k=1+\frac{20}{100}=1{,}2.
\]
Следовательно, подставлять в уравнение число 1,1 нельзя.

\textbf{Почему этот шаг неверен.} Если процентная ставка равна 20\%, то за год каждые 100 руб. превращаются в 120 руб., то есть вся сумма умножается на 1,2. Множитель 1,1 означает превращение 100 руб. в 110 руб., то есть увеличение только на 10\%. При использовании 1,1 вклад растёт медленнее, чем сказано в задаче; чтобы получить заданный конечный остаток, расчёт искусственно завышает первоначальный взнос. Отдельно отметим: $1{,}1^3=1{,}331$, но далее коэффициент при $S$ записан как 1,311. Это ещё одна несогласованность в вычислениях. Исправить лишь её недостаточно: основной ошибкой остаётся выбор ставки.

\idea\ В задаче на первоначальную сумму удобно выразить конечный вклад как сумму трёх частей: начальные деньги работают три года; пополнение перед вторым годом --- два года; пополнение перед третьим годом --- один год.

\textbf{Исправление от последнего правильного шага.} Сохраним составленное Вами уравнение, но возьмём правильный множитель:
\begin{align*}
173{,}28&=S(1{,}2)^3+10(1{,}2)^2+10(1{,}2),\\
173{,}28&=1{,}728S+14{,}4+12,\\
173{,}28&=1{,}728S+26{,}4,\\
1{,}728S&=173{,}28-26{,}4=146{,}88,\\
S&=\frac{146{,}88}{1{,}728}=85.
\end{align*}
Все суммы здесь даны в тысячах рублей.

\textbf{Полное правильное решение.} Пусть первоначальная сумма равна $S$ тыс. руб. За каждый год сумма умножается на 1,2.

\textit{Первый год.} Вклад $S$ приносит проценты; в конце года на счёте $1{,}2S$.

\textit{Начало второго года.} Добавляем 10 тыс. руб. Получаем $1{,}2S+10$. За второй год эта сумма возрастёт до
\[
(1{,}2S+10)\cdot1{,}2=1{,}44S+12.
\]
\textit{Начало третьего года.} Снова добавляем 10 тыс. руб. Сумма перед начислением процентов равна $1{,}44S+22$. После третьего года:
\[
(1{,}44S+22)\cdot1{,}2=1{,}728S+26{,}4.
\]
По условию это 173,28 тыс. руб. Приравниваем и решаем:
\[
1{,}728S+26{,}4=173{,}28
\quad\Longrightarrow\quad S=85.
\]
\textbf{Проверка.} При первоначальном вкладе 85 тыс. руб. после первого года будет $85\cdot1{,}2=102$ тыс. руб. После пополнения и второго года --- $(102+10)\cdot1{,}2=134{,}4$ тыс. руб. После второго пополнения и третьего года:
\[
(134{,}4+10)\cdot1{,}2=144{,}4\cdot1{,}2=173{,}28
\quad\text{тыс. руб.}
\]
Получена именно сумма из условия.

\result\ Первоначальная сумма вклада --- 85 тыс. руб., или 85\,000 руб. Структура Вашего уравнения была удачной; исправить нужно процентный множитель и связанные с ним вычисления.

\textbf{Задача на закрепление.} Вклад открыли на три года под 15\% годовых. В начале второго и третьего годов добавляли по 5 тыс. руб. В конце третьего года на счёте было 103\,615 руб. Найдите первоначальную сумму вклада в рублях.

\vfill
\begin{center}\footnotesize Лёвин Артём Александрович эксклюзивно для Анастасии Павловой\end{center}
\end{document}
